\documentclass[final,5p,times,twocolumn]{elsarticle}

\usepackage{amsmath,amssymb,amsfonts}
\usepackage{graphicx}
\usepackage{booktabs}
\usepackage{multirow}
\usepackage[hidelinks,hypertexnames=false]{hyperref}
\usepackage{algorithm}
\usepackage{algpseudocode}
\usepackage{float}
\usepackage{dblfloatfix}
\usepackage{microtype}
\usepackage{xcolor}
\usepackage{enumitem}
\usepackage{needspace}
\usepackage{capt-of}

\graphicspath{{./}{results/plots/}{workflowDiagrams/}}

% Balanced two-column float layout for diagrams and tables.
\setcounter{topnumber}{3}
\setcounter{dbltopnumber}{3}
\setcounter{totalnumber}{6}
\renewcommand{\topfraction}{0.92}
\renewcommand{\bottomfraction}{0.80}
\renewcommand{\textfraction}{0.06}
\renewcommand{\floatpagefraction}{0.82}
\renewcommand{\dbltopfraction}{0.92}
\renewcommand{\dblfloatpagefraction}{0.82}
\setlength{\textfloatsep}{11pt plus 2pt minus 2pt}
\setlength{\floatsep}{9pt plus 2pt minus 2pt}
\setlength{\intextsep}{9pt plus 2pt minus 2pt}
\hypersetup{hidelinks}

% Use supplied experimental plots when available. The original five plot files
% were not included with the source, so compact markers keep the draft compilable.
\newcommand{\resultplot}[1]{%
  \IfFileExists{results/plots/#1}{%
    \includegraphics[width=\columnwidth]{results/plots/#1}%
  }{%
    \begingroup
    \setlength{\fboxsep}{6pt}%
    \noindent\fbox{%
      \begin{minipage}[c][0.78in][c]{0.88\columnwidth}
        \centering\small\color{gray!75!black}
        Original result plot not supplied\\[3pt]
        \texttt{\detokenize{#1}}
      \end{minipage}%
    }%
    \endgroup
  }%
}

\journal{Future Generation Computer Systems}

\begin{document}

\begin{frontmatter}

\title{GranuWarm: Adaptive Dependency-Aware Pre-warming for Serverless Functions Under Non-Stationary Workloads Using Fuzzy-Controlled Constrained Particle Swarm Optimization}

\author[nsut]{Harshit Arya\corref{cor1}}
\ead{harshit.arya.ug24@nsut.ac.in}

\author[nsut]{Ayush Rawat}
\ead{ayush.rawat.ug24@nsut.ac.in}

\author[nsut]{Bhavya Yadav}
\ead{bhavya.yadav.ug24@nsut.ac.in}

\author[nsut]{Tijil Gupta}
\ead{tijil.gupta.ug24@nsut.ac.in}

\cortext[cor1]{Corresponding author}
\affiliation[nsut]{organization={Netaji Subhas University of Technology (NSUT)},
            city={New Delhi},
            country={India}}

\begin{abstract}
Serverless computing delivers elastic, event-driven execution of cloud functions without explicit infrastructure management, yet suffers from a persistent performance bottleneck: the \emph{cold start}---latency incurred when a container must be provisioned, initialized, and loaded with dependencies before execution begins. Cold-start mitigation is fundamentally a constrained optimization problem: pre-warming containers reduces latency but wastes resources on idle capacity, while aggressive reclamation conserves resources but increases user-facing latency and SLA violations. Existing approaches address fragments of this problem---static keep-alive policies, workload-prediction-based provisioning, isolated fuzzy or metaheuristic controllers, and dependency-aware scheduling---but to our knowledge, no framework jointly addresses reasoning under workload uncertainty, exploiting inter-function dependency structure for sub-container-level resource sharing, and solving the resulting constrained multi-function allocation problem under hard resource budgets.

This paper presents \textbf{GranuWarm}, a closed-loop adaptive prewarming framework that integrates three complementary mechanisms, each with a distinct architectural role. A \textbf{Mamdani-type fuzzy inference system} maps real-time workload signals---invocation intensity, traffic trend, cold-start pressure, and resource utilization---to per-function warm-capacity recommendations, enabling uncertainty-aware reasoning without requiring accurate long-horizon workload predictions. A \textbf{dependency granularity tree} decomposes the function fleet's software stacks into hierarchical, shareable layers (runtimes, libraries, function-specific code), transforming the decision space from coarse per-function container counts to fine-grained resource units that exploit inter-function overlap. A \textbf{constrained Particle Swarm Optimizer (PSO)} then searches for the optimal allocation of warm resources across granularity-tree nodes, minimizing a composite objective of cold-start latency, cold-start frequency, idle-resource cost, and SLA-violation rate, subject to platform memory and capacity constraints. GranuWarm employs a dual-timescale architecture: the fuzzy controller runs on a fast loop (10~s) for rapid reactive adaptation, while the PSO optimizer executes on a slower loop (60~s) for periodic global re-optimization. We evaluate GranuWarm using a trace-driven simulation environment calibrated against heterogeneous serverless microbenchmarks across five canonical workload regimes---stable, gradual-growth, bursty, periodic, and non-stationary---and conduct a systematic six-configuration ablation study with $N = 10$ paired trials across multiple random seeds. Experimental results demonstrate that GranuWarm achieves a superior Pareto frontier on the latency--resource tradeoff: it attains a P95 tail latency of 43.24~ms and an SLA violation rate of 0.90\% with an average memory footprint of 280.49~MB, reducing P95 latency by 99.7\% compared to scale-to-zero (15,015.68~ms) and by 98.2\% compared to static monolithic fixed prewarming (2,457.10~ms), while reducing memory usage by 95.2\% compared to monolithic full-container fixed prewarming (5,888.00~MB) and by over 89\% compared to non-granular dynamic baselines. Furthermore, the dependency granularity decomposition provides a 95.7\% reduction in composite fitness cost over non-granular fuzzy-PSO baselines, with two-tailed statistical significance ($p = 0.00195$) validated across all baseline comparisons.
\end{abstract}

\begin{keyword}
Serverless Computing \sep Cold-Start Mitigation \sep Fuzzy Inference System \sep Particle Swarm Optimization \sep Dependency-Aware Scheduling \sep Adaptive Resource Management
\end{keyword}

\end{frontmatter}

%% ====================================================================
%% 1. INTRODUCTION
%% ====================================================================
\section{Introduction}

\subsection{The Serverless Computing Paradigm}

Cloud computing has undergone a progressive shift from Infrastructure-as-a-Service (IaaS), through Platform-as-a-Service (PaaS), to Function-as-a-Service (FaaS), where the unit of deployment is a single stateless function~\cite{jonas2019cloud}. In the FaaS model---commercially realized by AWS Lambda, Google Cloud Functions, Azure Functions, and open-source platforms such as Knative~\cite{knative2024}, OpenFaaS~\cite{openfaas2024}, and Apache OpenWhisk~\cite{openwhisk2024}---developers deploy event-triggered functions without provisioning or managing servers. The platform handles container lifecycle management, automatic horizontal scaling, and crucially, \emph{scale-to-zero}: when a function receives no invocations for a configurable idle window, its containers are deallocated entirely, eliminating idle-resource costs.

This elasticity has driven rapid adoption. However, scale-to-zero introduces a well-documented performance challenge that remains a primary obstacle to serverless adoption for latency-sensitive applications: the \textbf{cold start}.

\subsection{Anatomy of a Cold Start}

When a function is invoked and a pre-initialized container is available (a \emph{warm start}), the request is dispatched directly, incurring only network and execution latency---typically single-digit milliseconds. When no warm container exists (a \emph{cold start}), the platform must execute a multi-stage initialization pipeline:

\begin{equation}
\begin{split}
T_{\text{cold}} &= T_{\text{sched}} + T_{\text{pull}} + T_{\text{runtime}} \\
&\quad + T_{\text{deps}} + T_{\text{init}} + T_{\text{exec}}
\end{split}
\end{equation}

where $T_{\text{sched}}$ is the scheduling and container-creation overhead, $T_{\text{pull}}$ is the container image pull time, $T_{\text{runtime}}$ is the language runtime startup (e.g., Python interpreter, Node.js V8), $T_{\text{deps}}$ is the loading and initialization of application-level dependencies (e.g., importing NumPy, loading ML weights), $T_{\text{init}}$ is function-specific initialization code, and $T_{\text{exec}}$ is the actual function execution time.

Empirical measurements across major cloud providers reveal that cold-start latencies range from approximately 100~ms for lightweight Node.js functions with minimal dependencies to over 10 seconds for functions with large dependency trees or ML model loading~\cite{mohan2019agile, wang2018peeking, bermbach2020using}. Shahrad et al.~\cite{shahrad2020serverless}, analyzing production traces from Azure Functions, found that cold starts affected up to 40\% of function invocations during low-traffic periods and contributed disproportionately to tail latency, with P99 latencies exceeding warm-start P99 by 10--50$\times$.

\subsection{The Latency--Resource Tradeoff}

Cold-start mitigation is fundamentally a \textbf{constrained multi-objective optimization} between two competing goals:

\textbf{Objective 1 --- Minimize cold-start impact and latency:}
\begin{equation}
\begin{split}
\min_{\mathbf{x} \in \mathcal{X}} \quad & \alpha_1 \cdot \mathbb{E}[T_{\text{cold}}(\mathbf{x})] + \alpha_2 \cdot P(\text{cold start} \mid \mathbf{x}) \\
& + \alpha_3 \cdot P(\text{SLA violation} \mid \mathbf{x})
\end{split}
\end{equation}

\textbf{Objective 2 --- Minimize resource waste:}
\begin{equation}
\min_{\mathbf{x} \in \mathcal{X}} \beta_1 \cdot M_{\text{idle}}(\mathbf{x}) + \beta_2 \cdot C_{\text{idle}}(\mathbf{x})
\end{equation}

where $M_{\text{idle}}$ and $C_{\text{idle}}$ represent idle memory and CPU consumption from pre-warmed but unused capacity.

These objectives are in direct tension. Aggressive prewarming minimizes cold starts but increases resource waste. Aggressive reclamation conserves resources but spikes tail latency. This tradeoff is compounded by burstiness, non-stationarity, runtime heterogeneity, and hard platform resource budgets. The central question is therefore:

\begin{quote}
\emph{Given observed workload dynamics, finite platform resources, and per-function SLA requirements, what is the optimal allocation of warm capacity across the function fleet---at what granularity---and how should this allocation adapt over time?}
\end{quote}

\subsection{Limitations of Existing Approaches and Prior Work}

Existing work addresses fragments of this challenge:
\begin{itemize}[leftmargin=*]
    \item \textbf{Static keep-alive policies:} Fixed keep-alive timers (e.g., 10-minute idle window) or static container allocations waste substantial resources during low-traffic periods and provide no adaptive protection against traffic surges~\cite{akkus2018sand, oakes2018sock}.
    \item \textbf{Prediction-based prewarming:} Predictive models (ARIMA, LSTM, regression) forecast upcoming invocations~\cite{suresh2021servermore, shahrad2019architectural, daraghmeh2024optimizing, lin2021modeling}. However, they are sensitive to forecast errors during sudden regime shifts.
    \item \textbf{Sub-container layer sharing:} Recent systems like RainbowCake~\cite{rainbowcake2024}, Pagurus~\cite{pagurus2022}, and FaasCache~\cite{faascache2021} highlight the value of fine-grained, layer-wise container caching and runtime sharing. Hanaforoosh et al.~\cite{hanaforoosh2025granularity} proposed a granularity-tree scheduler for dependency reuse. However, these systems focus on the runtime mechanism rather than an adaptive, closed-loop multi-function control policy under hard platform memory budgets.
    \item \textbf{Metaheuristic optimization:} Standalone genetic algorithms or PSO have been applied to container placement~\cite{buyya2023pso, zhan2022survey}, but they lack uncertainty-aware feedback controllers.
\end{itemize}

As summarized in Table~\ref{tab:gap_analysis}, to our knowledge, no existing framework simultaneously integrates workload-uncertainty-aware fuzzy control, hierarchical dependency-granularity decomposition, and constrained metaheuristic optimization in a closed control loop.

\begin{table}[htbp]
\centering
\caption{Coverage of existing approaches across problem dimensions.}
\label{tab:gap_analysis}
\resizebox{\columnwidth}{!}{%
\begin{tabular}{lccccc}
\toprule
Approach & \rotatebox{45}{Workload Unc.} & \rotatebox{45}{Multi-Func. Constr.} & \rotatebox{45}{Dep. Awareness} & \rotatebox{45}{Adaptive Ctrl.} & \rotatebox{45}{Online Opt.} \\
\midrule
Static policies~\cite{akkus2018sand, oakes2018sock} & \texttimes & \texttimes & \texttimes & \texttimes & \texttimes \\
Prediction-based~\cite{suresh2021servermore, daraghmeh2024optimizing} & Partial & \checkmark & \texttimes & \texttimes & \checkmark \\
Granularity-tree~\cite{hanaforoosh2025granularity} & \texttimes & \texttimes & \checkmark & \texttimes & \texttimes \\
Layer caching (RainbowCake)~\cite{rainbowcake2024} & \texttimes & Partial & \checkmark & \texttimes & \texttimes \\
PSO-based~\cite{buyya2023pso} & \texttimes & \checkmark & \texttimes & \texttimes & \checkmark \\
\textbf{GranuWarm (ours)} & \checkmark & \checkmark & \checkmark & \checkmark & \checkmark \\
\bottomrule
\end{tabular}%
}
\end{table}

\subsection{Proposed Approach: GranuWarm}

We propose \textbf{GranuWarm}, an adaptive closed-loop prewarming framework structured around three coordinated mechanisms:
\begin{enumerate}
    \item \textbf{Mamdani Fuzzy Inference System} $\rightarrow$ Operates on a fast control loop (10~s) to evaluate real-time signals (invocation intensity, traffic trend, cold-start pressure, memory utilization) and infer desired warm capacities without requiring rigid predictive models.
    \item \textbf{Dependency Granularity Tree} $\rightarrow$ Decomposes monolithic function environments into hierarchical, shareable nodes (language runtimes, library bundles, leaf code), mapping function demands into shared-layer capacities.
    \item \textbf{Constrained PSO Optimizer} $\rightarrow$ Executes on a slower loop (60~s) to find the globally optimal, budget-compliant allocation vector across granularity-tree nodes, seeded with the fuzzy recommendation.
\end{enumerate}

\subsection{Contributions}

The principal contributions of this work are:
\begin{enumerate}
    \item \textbf{Granularity-aware formulation:} A formal model mapping serverless function execution to hierarchical dependency trees, explicitly defining partial warm-start latency reductions.
    \item \textbf{Dual-timescale closed-loop architecture:} A control architecture coupling rapid fuzzy inference (10~s) with constrained discrete PSO (60~s), providing fast adaptation while preventing optimizer overhead.
    \item \textbf{Systematic empirical evaluation and ablation:} A comprehensive evaluation across five canonical workload regimes with $N = 10$ paired trials, demonstrating that GranuWarm achieves an optimal latency--resource tradeoff and that dependency-granularity tree modeling contributes the single largest reduction in resource footprint and tail latency.
\end{enumerate}

%% ====================================================================
%% 2. RELATED WORK
%% ====================================================================
\section{Related Work}

\subsection{Cold-Start Characterization and Keep-Alive Strategies}

Mohan et al.~\cite{mohan2019agile} and Wang et al.~\cite{wang2018peeking} established that cold-start overheads are dominated by dependency loading and runtime initialization rather than container provisioning alone. Shahrad et al.~\cite{shahrad2020serverless} characterized production Azure Functions traces, showing that invocation inter-arrival times span several orders of magnitude and that fixed keep-alive timers result in severe resource inefficiency. Suresh et al.~\cite{suresh2021servermore} and Daraghmeh et al.~\cite{daraghmeh2024optimizing} evaluated predictive windowing and regression models for function invocations.

\subsection{Layer-Aware and Fine-Grained Container Sharing}

To avoid deploying duplicate dependencies across containers, recent systems have explored sub-container resource sharing. RainbowCake~\cite{rainbowcake2024} implemented layer-wise container caching and sharing, demonstrating significant memory savings. Pagurus~\cite{pagurus2022} and FaasCache~\cite{faascache2021} designed container sharing and checkpointing primitives. Hanaforoosh et al.~\cite{hanaforoosh2025granularity} introduced granularity trees to represent function packages as hierarchical DAGs. GranuWarm builds on the granularity-tree concept but introduces a closed-loop adaptive control mechanism that dynamically optimizes warm capacity across shared layers under platform memory constraints.

\subsection{Fuzzy Logic and Metaheuristics in Cloud Systems}

Fuzzy logic controllers have been utilized for autoscaling and resource provisioning under uncertainty~\cite{jamshidi2014autonomic, mamdani1975experiment}. Concurrently, metaheuristics like PSO~\cite{kennedy1995pso, buyya2023pso} have proven effective for combinatorial scheduling problems. GranuWarm bridges these paradigms by using fuzzy inference to seed and guide a constrained discrete PSO over dependency-tree nodes.

%% ====================================================================
%% 3. PROBLEM FORMULATION
%% ====================================================================
\section{Problem Formulation}

\subsection{System Model}

Consider a serverless platform hosting a set $F = \{f_1, f_2, \ldots, f_N\}$ of $N$ functions. Each function $f_i$ is defined by:
\begin{itemize}
    \item $m_i$: memory required per container instance (MB)
    \item $\tau_i^{\text{cold}}$: cold-start latency when initialized from scratch (ms)
    \item $\tau_i^{\text{warm}}$: warm-start execution latency (ms)
    \item $\lambda_i(t)$: invocation arrival rate at time $t$ (req/s)
    \item $\delta_i$: SLA latency threshold (ms)
\end{itemize}
The platform enforces a hard total memory budget $M_{\text{budget}}$ allocated for pre-warmed capacity.

\subsection{Dependency Granularity Model and Partial Warm Starts}

Each function $f_i$ has a dependency stack $D_i$. Merging common dependencies across functions yields a \textbf{granularity tree} $G = (V, E)$, where $V = \{v_1, v_2, \ldots, v_K\}$ represents distinct dependency layers ($K$ nodes total), ordered from root runtimes to application-specific leaves. Each node $v_k$ has an associated memory footprint $m_{v_k}$ and isolated initialization time $\tau_{v_k}$.

Let $V(f_i) \subseteq V$ denote the path of granularity-tree nodes required by function $f_i$. An allocation vector is $\mathbf{x} = [x_1, x_2, \ldots, x_K] \in \mathbb{Z}_{\geq 0}^K$, where $x_k$ is the number of warm instances maintained for node $v_k$.

\textbf{Partial Warm Start Latency:} When function $f_i$ is invoked, any node $v_k \in V(f_i)$ that has available warm capacity ($x_k \geq 1$) avoids its initialization latency. The resulting startup latency is:
\begin{equation}
\tau_i(\mathbf{x}) = \tau_i^{\text{warm}} + \sum_{v_k \in V(f_i), \, x_k = 0} \tau_{v_k}
\label{eq:partial_warm}
\end{equation}
As formulated in Equation~\eqref{eq:partial_warm}, if all nodes in $V(f_i)$ are warm, $\tau_i(\mathbf{x}) = \tau_i^{\text{warm}}$. If no nodes are warm, $\tau_i(\mathbf{x}) = \tau_i^{\text{warm}} + \sum_{v_k \in V(f_i)} \tau_{v_k} = \tau_i^{\text{cold}}$.

\subsection{Performance Metrics and Idle Accounting}

The effective warm instance availability for function $f_i$ is $w_i(\mathbf{x}) = \min_{v_k \in V(f_i)} x_k$.

The \textbf{cold-start rate} for function $f_i$ over monitoring window $T_{\text{window}}$ is:
\begin{equation}
C_i(\mathbf{x}) = \max\left(0, 1 - \frac{w_i(\mathbf{x})}{\max(1, \lceil \lambda_i \cdot T_{\text{window}} \rceil)}\right)
\end{equation}

The \textbf{expected latency} for function $f_i$ is:
\begin{equation}
L_i(\mathbf{x}) = C_i(\mathbf{x}) \cdot \tau_i(\mathbf{x}) + (1 - C_i(\mathbf{x})) \cdot \tau_i^{\text{warm}}
\end{equation}

Let $u_k$ denote the number of active instances of node $v_k$ currently processing requests. The \textbf{idle memory cost} is:
\begin{equation}
M_{\text{idle}}(\mathbf{x}) = \sum_{k=1}^{K} m_{v_k} \cdot \max(0, x_k - u_k)
\end{equation}

The \textbf{SLA violation indicator} is:
\begin{equation}
S_i(\mathbf{x}) = \mathbb{I}[L_i(\mathbf{x}) > \delta_i]
\end{equation}

\subsection{Optimization Problem}

The prewarming allocation problem is formulated as:
\begin{equation}
\begin{split}
\min_{\mathbf{x}} \quad J(\mathbf{x}) &= w_L \frac{\bar{L}(\mathbf{x})}{L_{\max}} + w_C \bar{C}(\mathbf{x}) \\
&\quad + w_M \frac{M_{\text{idle}}(\mathbf{x})}{M_{\text{budget}}} + w_S \bar{S}(\mathbf{x})
\end{split}
\end{equation}
subject to:
\begin{align}
\sum_{k=1}^{K} m_{v_k} \cdot x_k &\leq M_{\text{budget}} \quad \text{(hard memory budget)} \\
x_k^{\min} \leq x_k &\leq x_k^{\max} \quad \forall k \in \{1, \ldots, K\} \quad \text{(bounds)} \\
x_k &\in \mathbb{Z}_{\geq 0} \quad \forall k \quad \text{(integrality)}
\end{align}
where $\bar{L}(\mathbf{x}) = \frac{1}{N}\sum_i L_i(\mathbf{x})$, $\bar{C}(\mathbf{x}) = \frac{1}{N}\sum_i C_i(\mathbf{x})$, $\bar{S}(\mathbf{x}) = \frac{1}{N}\sum_i S_i(\mathbf{x})$, and $L_{\max}$ is the maximum cold-start latency across all functions. Default weights: $w_L = 0.35, w_C = 0.25, w_M = 0.25, w_S = 0.15$. Bounds are set to $x_k^{\min} = 0, x_k^{\max} = 10$, and $M_{\text{budget}} = 4096$~MB.

%% ====================================================================
%% 4. PROPOSED FRAMEWORK
%% ====================================================================
\section{Proposed Framework: GranuWarm}

\subsection{System Architecture Overview}

GranuWarm operates as a closed-loop controller with six functional components:
\begin{enumerate}
    \item \textbf{Workload Monitor:} Gathers invocation rates, trends, cold-start counts, and memory usage.
    \item \textbf{Fuzzy Inference System:} Maps real-time workload signals to per-function capacity recommendations on a fast 10~s loop.
    \item \textbf{Dependency Granularity Tree:} Projects function recommendations onto shared nodes.
    \item \textbf{Constrained PSO Optimizer:} Periodically re-optimizes global allocations on a 60~s loop under hard memory constraints.
    \item \textbf{Actuation Engine:} Applies allocation diffs with hysteresis.
    \item \textbf{Feedback Loop:} Feeds observed latencies and memory usage back into the monitor.
\end{enumerate}

\subsection{Component 1: Workload Monitor}

Every $\Delta t_{\text{fast}} = 10$~s, the monitor extracts four normalized signals:
\begin{itemize}[leftmargin=*]
    \item \textbf{Invocation Intensity} $R_i(t) \in [0, 1]$: Request rate normalized by the function's historical peak.
    \item \textbf{Traffic Trend} $T_i(t) \in [-1, 1]$: First-order difference $R_i(t) - R_i(t - \Delta t)$.
    \item \textbf{Cold-Start Pressure} $C_i(t) \in [0, 1]$: Ratio of cold starts to total requests in sliding window $W = 60$~s.
    \item \textbf{Resource Pressure} $M(t) \in [0, 1]$: Ratio of total allocated memory to $M_{\text{budget}}$.
\end{itemize}

\subsection{Component 2: Fuzzy Inference System}

The fuzzy system uses standard triangular and trapezoidal membership functions:
\begin{itemize}[leftmargin=*]
    \item $R$: LOW $[0, 0, 0.15, 0.35]$, MEDIUM $[0.2, 0.5, 0.8]$, HIGH $[0.65, 0.85, 1, 1]$.
    \item $T$: FALLING $[-1, -1, -0.5, -0.1]$, STABLE $[-0.2, 0, 0.2]$, RISING $[0.1, 0.4, 0.7]$, SURGING $[0.5, 0.8, 1, 1]$.
    \item $C$: LOW $[0, 0, 0.1, 0.3]$, MEDIUM $[0.15, 0.4, 0.65]$, HIGH $[0.5, 0.7, 1, 1]$.
    \item $M$: LOW $[0, 0, 0.3, 0.5]$, MEDIUM $[0.35, 0.55, 0.75]$, HIGH $[0.65, 0.8, 1, 1]$.
    \item Output $P$ (Desired Warm Capacity): VERY\_LOW $[0, 0, 0.05, 0.2]$, LOW $[0.1, 0.25, 0.4]$, MEDIUM $[0.3, 0.5, 0.7]$, HIGH $[0.6, 0.75, 0.9]$, VERY\_HIGH $[0.8, 0.95, 1, 1]$.
\end{itemize}

Table~\ref{tab:fuzzy_rules} defines the 27 operational rules. Rule evaluation uses min-operator for conjunction, max-operator for aggregation, and centroid defuzzification. Defuzzified output $P_i^{\text{fuzzy}} \in [0, 1]$ scales to:
\begin{equation}
P_i^{\text{desired}} = x_i^{\min} + P_i^{\text{fuzzy}} \cdot (x_i^{\max} - x_i^{\min})
\end{equation}

\begin{table}[H]
\centering
\caption{Complete fuzzy rule base (27 rules).}
\label{tab:fuzzy_rules}
\resizebox{\columnwidth}{!}{%
\begin{tabular}{clllll}
\toprule
\# & Invocation ($R$) & Trend ($T$) & Cold-Start ($C$) & Memory ($M$) & $\rightarrow$ Warm Cap.~($P$) \\
\midrule
1 & LOW & FALLING & LOW & --- & VERY\_LOW \\
2 & LOW & STABLE & LOW & --- & VERY\_LOW \\
3 & LOW & STABLE & MEDIUM & LOW & LOW \\
4 & LOW & STABLE & MEDIUM & MEDIUM & VERY\_LOW \\
5 & LOW & RISING & LOW & LOW & LOW \\
6 & LOW & RISING & MEDIUM & LOW & MEDIUM \\
7 & LOW & RISING & HIGH & LOW & MEDIUM \\
8 & LOW & SURGING & --- & LOW & MEDIUM \\
9 & LOW & SURGING & --- & HIGH & LOW \\
10 & MEDIUM & FALLING & LOW & --- & LOW \\
11 & MEDIUM & STABLE & LOW & LOW & MEDIUM \\
12 & MEDIUM & STABLE & LOW & HIGH & LOW \\
13 & MEDIUM & STABLE & MEDIUM & LOW & MEDIUM \\
14 & MEDIUM & STABLE & HIGH & LOW & HIGH \\
15 & MEDIUM & RISING & LOW & LOW & MEDIUM \\
16 & MEDIUM & RISING & MEDIUM & LOW & HIGH \\
17 & MEDIUM & RISING & HIGH & LOW & HIGH \\
18 & MEDIUM & SURGING & --- & LOW & HIGH \\
19 & MEDIUM & SURGING & --- & HIGH & MEDIUM \\
20 & HIGH & FALLING & LOW & --- & MEDIUM \\
21 & HIGH & STABLE & LOW & LOW & HIGH \\
22 & HIGH & STABLE & MEDIUM & LOW & HIGH \\
23 & HIGH & STABLE & HIGH & LOW & VERY\_HIGH \\
24 & HIGH & RISING & --- & LOW & VERY\_HIGH \\
25 & HIGH & RISING & --- & HIGH & HIGH \\
26 & HIGH & SURGING & --- & LOW & VERY\_HIGH \\
27 & HIGH & SURGING & --- & HIGH & HIGH \\
\bottomrule
\end{tabular}%
}
\end{table}

% Figure 1: compact one-column diagram immediately after the fuzzy component.
\begin{figure}[H]
    \centering
    \includegraphics[width=0.98\columnwidth,height=0.22\textheight,keepaspectratio]{figures/fuzzy_inference_pipeline.png}
    \caption{Mamdani fuzzy inference pipeline mapping real-time workload metrics to per-function warm-capacity recommendations.}
    \label{fig:fuzzy_pipeline}
\end{figure}

\subsection{Component 3: Dependency Granularity Tree}

Table~\ref{tab:gran_example} lists the complete set of $K = 20$ granularity-tree nodes constructed from our benchmark fleet.

\begin{table}[htbp]
\centering
\caption{Granularity tree node structure ($K=20$ nodes total).}
\label{tab:gran_example}
\resizebox{\columnwidth}{!}{%
\begin{tabular}{llrrl}
\toprule
Layer & Node ID & Mem (MB) & Init (ms) & Dependent Functions / Children \\
\midrule
$-$1 & System Root & 0 & 0 & System root \\
0 & Python 3.11 Runtime & 80 & 200 & F03--F09, F11 \\
0 & Node.js 20 Runtime & 40 & 100 & F01, F02, F10, F12 \\
0 & OS Native Shared C-Libs & 15 & 20 & Python and Node.js runtimes \\
1 & NumPy Core Bundle & 25 & 80 & F03, F04, F06, F08, F09 \\
1 & Pillow Image Bundle & 15 & 50 & F03, F04, F05 \\
1 & Pandas Data Bundle & 35 & 90 & F06, F08, F09 \\
1 & Shared Web Utilities & 10 & 30 & F02, F10 \\
2 & F01 helloWorld Leaf & 1 & 10 & F01 \\
2 & F02 jsonTransform Leaf & 5 & 20 & F02 \\
2 & F03 imageResize Leaf & 5 & 30 & F03 \\
2 & F04 imageFilter Leaf & 60 & 120 & F04 \\
2 & F05 thumbnailGen Leaf & 3 & 20 & F05 \\
2 & F06 mlInference Leaf & 400 & 800 & F06 \\
2 & F07 sentimentAnalysis Leaf & 1800 & 3500 & F07 \\
2 & F08 dataAggregate Leaf & 5 & 30 & F08 \\
2 & F09 csvParser Leaf & 3 & 20 & F09 \\
2 & F10 authService Leaf & 8 & 40 & F10 \\
2 & F11 orderProcess Leaf & 30 & 150 & F11 \\
2 & F12 notification Leaf & 15 & 60 & F12 \\
\bottomrule
\end{tabular}%
}
\end{table}

\textbf{Projection of Fuzzy Recommendations:} For each node $v_k$, let $F(v_k)$ denote the set of functions depending on $v_k$. The projected demand on node $v_k$ is:
\begin{equation}
P_k^{\text{tree}} = \max_{f_i \in F(v_k)} P_i^{\text{desired}}
\end{equation}

% Figure 2: the hierarchical tree is compact enough for one column and follows
% the full Component 3 explanation, including the projection equation.
\Needspace{0.31\textheight}
\begin{figure}[H]
    \centering
    \includegraphics[width=0.88\columnwidth,height=0.30\textheight,keepaspectratio]{figures/dependency_granularity_tree.png}
    \caption{Dependency granularity tree showing hierarchical decomposition of function software stacks into runtime, shared-library, and function-specific layers.}
    \label{fig:granularity_tree}
\end{figure}

\subsection{Component 4: Constrained PSO Optimizer}

A swarm of $S = 30$ particles searches the integer allocation vector $\mathbf{x} \in \mathbb{Z}_{\geq 0}^K$.
\begin{itemize}[leftmargin=*]
    \item \textbf{Informed Seeding:} Particle 1 is initialized with the projected fuzzy recommendation $P_k^{\text{tree}}$; Particle 2 with the current active allocation $\mathbf{x}_{\text{current}}$; Particles 3--5 with scaled variants ($0.5\times, 0.75\times, 1.25\times$); and Particles 6--30 with random perturbations around $P_k^{\text{tree}}$.
    \item \textbf{Budget Constraint and Repair:} Candidates violating the memory budget ($\sum_{k=1}^K m_{v_k} x_k > M_{\text{budget}}$) are penalized using Deb's parameter-less constraint penalty method~\cite{deb2000constraint}. Furthermore, candidate solutions undergo a deterministic greedy budget repair: the algorithm iteratively decrements $x_j$ for node $j = \arg\max_{k: x_k > x_k^{\min}} (m_{v_k} / \tau_{v_k})$---pruning nodes with the highest memory cost per millisecond of initialization benefit---until the hard platform budget $M_{\text{budget}}$ is strictly satisfied.
    \item \textbf{Dynamics:} Standard discrete PSO equations update positions over $I_{\max} = 50$ iterations with inertia $\omega$ decaying from 0.7 to 0.4 and cognitive/social coefficients $c_1 = c_2 = 1.5$.
\end{itemize}

% Figure 3: single-column workflow placed directly after Component 4.
% Keeping this diagram in one column preserves its local relationship with
% the PSO description without pushing the following subsection away.
\Needspace{0.22\textheight}
\begin{figure}[H]
    \centering
    \includegraphics[width=0.98\columnwidth]{figures/constrained_pso_workflow.png}
    \caption{Constrained PSO workflow using fuzzy recommendations, the granularity-tree structure, and resource constraints to determine warm-unit allocations.}
    \label{fig:pso_workflow}
\end{figure}
\vspace{8pt}

% Figure 4: standard double-column top float. This avoids the forced page
% break caused by a manual \twocolumn[...] block and lets the text flow on.
\begin{figure*}[!t]
    \centering
    \includegraphics[width=0.98\textwidth,height=0.27\textheight,keepaspectratio]{figures/dual_timescale_control_loop.png}
    \caption{GranuWarm dual-timescale control loop: fast fuzzy adaptation every 10~s and constrained PSO re-optimization every 60~s.}
    \label{fig:dual_timescale}
\end{figure*}

\Needspace{0.22\textheight}
\subsection{Dual-Timescale Control Loop}

Algorithm~\ref{alg:control_loop} details the dual-timescale coordination. To ensure that PSO optimization is preserved throughout the 60~s window, the fast loop applies bounded fuzzy adjustments $\Delta \mathbf{x}_{\text{fuzzy}}$ to the base PSO allocation $\mathbf{x}_{\text{pso}}^*$, rather than overwriting it.

\begin{algorithm}[H]
\caption{GranuWarm Dual-Timescale Control Loop}
\label{alg:control_loop}
\begin{algorithmic}[1]
\Require Function set $F$, Granularity tree $G$, Budget $M_{\text{budget}}$
\Require Fast interval $\Delta t_{\text{fast}} = 10$~s, Slow interval $\Delta t_{\text{slow}} = 60$~s
\Ensure Continuous adaptive allocation $\mathbf{x}$
\State $\mathbf{x}_{\text{pso}}^* \leftarrow \text{InitialUniformAllocation}(M_{\text{budget}})$
\State $\mathbf{x} \leftarrow \mathbf{x}_{\text{pso}}^*$; $t_{\text{last\_pso}} \leftarrow \text{current\_time}()$
\Loop
    \State $\text{metrics} \leftarrow \text{WorkloadMonitor.collect}()$
    \For{each function $f_i \in F$}
        \State $P_i \leftarrow \text{FuzzyInfer}(R_i, T_i, C_i, M)$
    \EndFor
    \State $P_{\text{tree}} \leftarrow \text{GranularityTree.project}(\{P_1, \ldots, P_N\})$
    \If{$\text{current\_time}() - t_{\text{last\_pso}} \geq \Delta t_{\text{slow}}$}
        \State $\mathbf{x}_{\text{pso}}^* \leftarrow \text{PSO.optimize}(P_{\text{tree}}, G, M_{\text{budget}}, \text{metrics})$
        \State $\mathbf{x} \leftarrow \mathbf{x}_{\text{pso}}^*$
        \State $t_{\text{last\_pso}} \leftarrow \text{current\_time}()$
    \Else
        \State $\Delta \mathbf{x} \leftarrow \text{BoundedDelta}(P_{\text{tree}}, \mathbf{x}_{\text{pso}}^*, \text{metrics})$
        \State $\mathbf{x} \leftarrow \text{RepairBudget}(\mathbf{x}_{\text{pso}}^* + \Delta \mathbf{x}, M_{\text{budget}})$
    \EndIf
    \State $\text{ActuationEngine.apply}(\mathbf{x})$
    \State \textbf{WAIT}($\Delta t_{\text{fast}}$)
\EndLoop
\end{algorithmic}
\end{algorithm}

%% ====================================================================
%% 5. EXPERIMENTAL METHODOLOGY
%% ====================================================================
\Needspace{0.22\textheight}
\section{Experimental Methodology}

\subsection{Benchmark Functions and Workload Regimes}

We evaluate on 12 heterogeneous benchmark functions (Table~\ref{tab:benchmark_functions}), spanning lightweight web handlers, image manipulation, machine learning inference, and data processing.



\begin{table*}[htbp]
\centering
\caption{Benchmark function characteristics and dependency profiles.}
\label{tab:benchmark_functions}
\begin{tabular}{llllrrl}
\toprule
ID & Function & Runtime & Key Dependencies & Mem (MB) & Cold-Start (ms) & Category \\
\midrule
F01 & helloWorld & Node.js 20 & Minimal & 128 & 150 & Lightweight \\
F02 & jsonTransform & Node.js 20 & lodash & 128 & 200 & Lightweight \\
F03 & imageResize & Python 3.11 & Pillow, numpy & 512 & 1,200 & Image \\
F04 & imageFilter & Python 3.11 & Pillow, numpy, scipy & 512 & 1,400 & Image \\
F05 & thumbnailGen & Python 3.11 & Pillow, numpy & 256 & 900 & Image \\
F06 & mlInference & Python 3.11 & scikit-learn, numpy, pandas & 1,024 & 3,500 & ML \\
F07 & sentimentAnalysis & Python 3.11 & transformers, torch & 2,048 & 8,000 & ML \\
F08 & dataAggregate & Python 3.11 & pandas, numpy & 512 & 1,500 & Data \\
F09 & csvParser & Python 3.11 & pandas, numpy & 256 & 1,000 & Data \\
F10 & authService & Node.js 20 & jsonwebtoken, bcrypt & 256 & 300 & API \\
F11 & orderProcess & Python 3.11 & requests, sqlalchemy & 256 & 800 & API \\
F12 & notification & Node.js 20 & firebase-admin & 256 & 350 & API \\
\bottomrule
\end{tabular}
\end{table*}

Note that Table~\ref{tab:benchmark_functions} reports monolithic container memory allocations and full cold-start latencies when provisioned from scratch without layer sharing, whereas Table~\ref{tab:gran_example} details the decomposed incremental layer footprints and isolated startup overheads utilized by the granularity tree model.

Experiments evaluate five canonical workload regimes (Table~\ref{tab:workloads}) for 30 minutes each:
\begin{itemize}[leftmargin=*]
    \item \textbf{W1 (Stable):} Constant Poisson arrival rates.
    \item \textbf{W2 (Growth):} Linearly increasing arrival rates ($+5\%$/min).
    \item \textbf{W3 (Bursty):} Periodic $10\times$ traffic bursts every 5 minutes.
    \item \textbf{W4 (Periodic):} Sinusoidal diurnal fluctuation with period $T=15$~min.
    \item \textbf{W5 (Non-Stationary):} Sudden regime shift at $t=T/2$ with base rate doubling and popularity inversion.
\end{itemize}

\begin{table}[htbp]
\centering
\caption{Workload regime characteristics.}
\label{tab:workloads}
\resizebox{\columnwidth}{!}{%
\begin{tabular}{llp{5.5cm}}
\toprule
Regime & Description & Invocation Pattern \\
\midrule
W1: Stable & Constant-rate Poisson arrivals & $\lambda_i(t) = \bar{\lambda}_i \;\forall t$ \\
W2: Growth & Linearly increasing rate & $\lambda_i(t) = \bar{\lambda}_i \cdot (1 + 0.05t)$ \\
W3: Bursty & Periodic traffic spikes (10$\times$) & $\lambda_i(t) = \bar{\lambda}_i + 9\bar{\lambda}_i \cdot \mathbf{1}[t \in \text{burst}]$ \\
W4: Periodic & Sinusoidal diurnal pattern & $\lambda_i(t) = \bar{\lambda}_i \cdot (1 + 0.8\sin(2\pi t / T))$ \\
W5: Non-Stat. & Regime change at $t = T/2$ & Rate shift + popularity inversion \\
\bottomrule
\end{tabular}%
}
\end{table}

% Figure 5: workload-regime patterns follow the workload description and table.
\begin{figure*}[!t]
    \centering
    \includegraphics[width=0.98\textwidth,height=0.27\textheight,keepaspectratio]{figures/workload_regimes.png}
    \caption{Illustrative invocation-rate patterns for the five workload regimes (W1--W5).}
    \label{fig:workload_regimes}
\end{figure*}




\subsection{Experimental Configurations (Ablation Setup)}

We compare GranuWarm against five baseline configurations in a systematic ablation:
\begin{itemize}[leftmargin=*]
    \item \textbf{B1 (No Prewarming):} Default scale-to-zero without warm buffers.
    \item \textbf{B2 (Fixed Prewarming):} Static prewarming allocation provisioning one monolithic container per function across all 12 benchmark functions without sub-container layer sharing ($\sum_{i=1}^{12} m_i = 5,888$~MB). We also analyze its performance when capped under hard platform memory budgets.
    \item \textbf{B3 (Fuzzy-Only):} Reactive fuzzy inference without PSO or granularity tree.
    \item \textbf{B4 (PSO-Only):} Periodic PSO running on per-function search space.
    \item \textbf{B5 (Fuzzy + PSO, No Granularity):} Integrated fuzzy and PSO operating on per-function container units.
    \item \textbf{B6 (GranuWarm Full):} Complete framework with fuzzy control, dependency granularity tree, and constrained PSO.
\end{itemize}

Statistical significance is assessed via the non-parametric Wilcoxon signed-rank test ($\alpha = 0.05$) across an expanded matrix of $N = 10$ paired trials (5 workload regimes $\times$ 2 independent random seeds, Seeds 42 and 123) evaluated under standard two-tailed hypothesis testing.

%% ====================================================================
%% 6. RESULTS AND DISCUSSION
%% ====================================================================
\section{Results and Discussion}

\subsection{Overall Benchmark Performance}

Table~\ref{tab:overall_perf} reports the overall performance across all baseline configurations, computed from the complete experimental matrix across workloads W1--W5. Figures~\ref{fig:pareto} and \ref{fig:csr} plot the Pareto frontier and cold-start rates, while Figure~\ref{fig:temporal} visualizes temporal dynamics under bursty traffic. Figure~\ref{fig:ablation} summarizes the ablation fitness costs, and Figure~\ref{fig:live_validation} presents live Kubernetes validation results.

\begin{figure}[htbp]
    \centering
    \resultplot{figure_1_pareto.pdf}
    \caption{Latency vs. Memory Pareto Frontier for baseline configurations.}
    \label{fig:pareto}
\end{figure}

\begin{figure}[htbp]
    \centering
    \resultplot{figure_2_csr.pdf}
    \caption{Cold-Start Frequency Across Workload Regimes (W1--W5).}
    \label{fig:csr}
\end{figure}

\begin{figure}[htbp]
    \centering
    \resultplot{figure_3_temporal.pdf}
    \caption{Temporal Dynamics Under Bursty W3 Traffic: GranuWarm (B6) rapidly dampens latency spikes via dual-timescale adaptation, while static and non-granular baselines exhibit sustained queuing delays.}
    \label{fig:temporal}
\end{figure}

\begin{table}[htbp]
\centering
\caption{Overall Benchmark Performance Across Baselines (Mean Values)}
\label{tab:overall_perf}
\resizebox{\columnwidth}{!}{%
\begin{tabular}{lccccc}
\toprule
Config & CSR (\%) & P95 Latency (ms) & Avg Mem (MB) & SVR (\%) & Composite $J$ \\
\midrule
B1 & 0.19 & 15,015.68 & 3,225.60 & 9.24 & 1.69 \\
B2 & 4.21 & 2,457.10 & 5,888.00 & 12.45 & 1.05 \\
B3 & 2.88 & 7,641.40 & 3,123.20 & 9.62 & 1.81 \\
B4 & 0.03 & 15,017.62 & 2,845.91 & 13.37 & 2.13 \\
B5 & 0.07 & 15,015.65 & 2,794.73 & 9.13 & 1.64 \\
B6 & 0.54 & 43.24 & 280.49 & 0.90 & 0.07 \\
\bottomrule
\end{tabular}%
}
\end{table}

\textbf{Latency and SLA Violations:} GranuWarm (B6) achieves the lowest P95 tail latency (43.24~ms) and the lowest SLA violation rate (0.90\%) across all configurations. Compared to default scale-to-zero (B1: 15,015.68~ms P95, 9.24\% SVR), GranuWarm virtually eliminates tail latency timeouts, achieving a 99.7\% reduction in P95 latency. Against static fixed prewarming (B2: 2,457.10~ms P95, 12.45\% SVR), GranuWarm achieves a 98.2\% reduction in P95 latency and a 92.8\% reduction in SLA violations, because fixed monolithic allocations cannot dynamically scale up to absorb concurrency bursts.

\textbf{Memory Footprint:} GranuWarm maintains an average memory footprint of only 280.49~MB. In sharp contrast, non-granular fixed prewarming (B2) provisioning a single monolithic container for each of the 12 functions requires $\sum_{i=1}^{12} m_i = 5,888.00$~MB of memory, exceeding the platform budget ($M_{\text{budget}} = 4,096$~MB) by 43.8\% and consuming $21.0\times$ more memory than GranuWarm. If B2 were strictly constrained to operate under a memory budget matching GranuWarm ($\sim$280~MB), it could warm at most 1--2 small containers (e.g., F01 and F02 = 256~MB), forcing the remaining 10 functions to cold start and queue on every invocation, driving P95 latency to $>15,000$~ms. Non-granular dynamic baselines (B1, B3, B4, B5) similarly consume 2,794--3,225~MB because each container instance independently duplicates full runtime and dependency stacks. GranuWarm achieves a 95.2\% memory reduction compared to monolithic B2 and over 89\% compared to non-granular baselines by sharing base runtimes and library bundles.

\textbf{Cold-Start Rate vs.\ Tail Latency:} A prominent phenomenon observed in Table~\ref{tab:overall_perf} is that baselines B4 and B5 exhibit near-zero cold-start rates (0.03\% and 0.07\%, respectively), yet experience extreme P95 tail latencies (15,017.62~ms and 15,015.65~ms). Because non-granular baselines (B4/B5) allocate full monolithic containers, their memory footprint hits the 4,096 MB budget after pre-warming only 2--3 large functions (e.g., ML models). The remaining 9 functions receive zero warm capacity. While the few allocated functions yield 0\% cold starts (lowering overall fleet average CSR), unallocated functions encounter severe concurrency queuing delays that hit 15-second execution timeouts, causing extreme P95 tail latencies. Similarly, under default scale-to-zero (B1), cold starts occur primarily during initial container provisioning; however, without pre-warmed capacity buffering incoming bursts, subsequent concurrent requests encounter severe queuing behind containers undergoing initialization, hitting platform timeout thresholds ($\sim$15~s) and causing extreme tail latencies (15,015.68~ms P95) and high SLA violations (9.24\%) even while steady-state cold-start event fractions appear low (0.19\%). In contrast, GranuWarm eliminates timeout spikes by proactively warming shared base layers across the entire fleet.

\subsection{Performance Across Workload Regimes}

Tables~\ref{tab:p95_latency}, \ref{tab:csr_workload}, and \ref{tab:mem_workload} present performance breakdowns across the five workload regimes.

\begin{table}[htbp]
\centering
\caption{P95 Tail Latency (ms) by Workload Regime}
\label{tab:p95_latency}
\resizebox{\columnwidth}{!}{%
\begin{tabular}{lccccc}
\toprule
Config & W1 (Stable) & W2 (Growth) & W3 (Bursty) & W4 (Periodic) & W5 (Non-Stat.) \\
\midrule
B1 & 15,010.5 & 15,015.0 & 15,021.7 & 15,016.2 & 15,014.9 \\
B2 & 1,463.6 & 1,485.2 & 4,020.4 & 1,470.8 & 3,845.5 \\
B3 & 35.6 & 15,009.0 & 23,050.3 & 49.2 & 62.9 \\
B4 & 15,017.6 & 15,019.2 & 15,021.5 & 15,016.4 & 15,013.4 \\
B5 & 15,010.4 & 15,014.5 & 15,021.4 & 15,017.5 & 15,014.4 \\
B6 & 34.0 & 37.3 & 66.6 & 40.1 & 38.2 \\
\bottomrule
\end{tabular}%
}
\end{table}

\begin{table}[htbp]
\centering
\caption{Cold-Start Rate (\%) by Workload Regime}
\label{tab:csr_workload}
\resizebox{\columnwidth}{!}{%
\begin{tabular}{lccccc}
\toprule
Config & W1 (Stable) & W2 (Growth) & W3 (Bursty) & W4 (Periodic) & W5 (Non-Stat.) \\
\midrule
B1 & 0.16 & 0.00 & 0.76 & 0.04 & 0.01 \\
B2 & 0.82 & 1.15 & 14.85 & 1.04 & 3.20 \\
B3 & 0.00 & 0.00 & 14.04 & 0.00 & 0.38 \\
B4 & 0.05 & 0.03 & 0.01 & 0.02 & 0.03 \\
B5 & 0.00 & 0.01 & 0.19 & 0.02 & 0.13 \\
B6 & 0.60 & 0.37 & 1.22 & 0.31 & 0.22 \\
\bottomrule
\end{tabular}%
}
\end{table}

\begin{table}[htbp]
\centering
\caption{Average Memory Footprint (MB) by Workload Regime}
\label{tab:mem_workload}
\resizebox{\columnwidth}{!}{%
\begin{tabular}{lccccc}
\toprule
Config & W1 (Stable) & W2 (Growth) & W3 (Bursty) & W4 (Periodic) & W5 (Non-Stat.) \\
\midrule
B1 & 3,072.0 & 3,968.0 & 3,328.0 & 2,048.0 & 3,712.0 \\
B2 & 5,888.0 & 5,888.0 & 5,888.0 & 5,888.0 & 5,888.0 \\
B3 & 3,072.0 & 3,072.0 & 3,072.0 & 3,328.0 & 3,072.0 \\
B4 & 3,840.0 & 3,968.0 & 3,072.0 & 3,072.0 & 277.5 \\
B5 & 4,096.0 & 3,456.0 & 3,072.0 & 3,072.0 & 277.6 \\
B6 & 277.4 & 279.4 & 287.9 & 280.2 & 277.5 \\
\bottomrule
\end{tabular}%
}
\end{table}

Under sudden popularity inversions in W5, non-granular PSO (B4/B5) violates memory constraints due to mispredicted demand, triggering greedy budget repair that strips container allocations down to minimum bounds (277 MB), resulting in massive cold-start queues.

\textbf{Bursty Workload (W3):} Under W3 traffic spikes, fixed prewarming (B2) with static single-container capacity experiences significant queuing delays when concurrent arrivals exceed unit capacity (P95 of 4,020.4~ms), while reactive fuzzy-only prewarming (B3) surges to 23,050.3~ms because reactive scaling alone cannot anticipate sudden $10\times$ bursts without global budget rebalancing. In contrast, GranuWarm (B6) dampens burst latency spikes to 66.6~ms P95, demonstrating that fine-grained shared layers provide elastic buffer capacity without container bloat. Figure~\ref{fig:temporal} illustrates this temporal behavior, showing how GranuWarm's dual-timescale controller rapidly re-converges after each burst event.

\subsection{Addressing Research Questions}

\subsubsection{RQ1: Does Workload-Aware Fuzzy Control Outperform Static Policies?}
Fixed prewarming (B2) illustrates the fundamental failure of static monolithic provisioning: prewarming one instance of each function wastes 5,888~MB of memory ($21.0\times$ GranuWarm), yet still suffers elevated tail latency (2,457.10~ms P95) and SLA violations (12.45\%) during traffic bursts because fixed instances cannot elastically absorb concurrency surges. If constrained to a tight budget, B2 suffers massive cold starts. Reactive fuzzy-only control (B3) succeeds under stable and periodic regimes (W1: 35.6~ms P95, W4: 49.2~ms P95) where traffic evolves smoothly, but encounters severe degradation under bursts (W3: 23,050.3~ms P95). Combining fuzzy recommendations with global optimization and granularity awareness (B6) provides dynamic elasticity within 280.49~MB, reducing composite fitness cost $J$ by 93.3\% compared to B2 (0.07 vs.\ 1.05).

\subsubsection{RQ2: Does Constrained PSO Improve the Tradeoff?}
Comparing B4 (PSO-Only: $J = 2.13$) and B5 (Fuzzy+PSO: $J = 1.64$) against B3 demonstrates the utility of PSO. Seeding the PSO with fuzzy recommendations guides particles into viable allocation regions, yielding faster convergence and a 23.0\% improvement in composite cost over standalone PSO.

\begin{figure}[htbp]
    \centering
    \resultplot{figure_4_ablation.pdf}
    \caption{Ablation Study Performance Progression showing composite fitness cost $J$ across configurations B1--B6. Each bar represents the mean $J$ across all workload regimes, with error bars indicating standard deviation across $N=10$ paired trials.}
    \label{fig:ablation}
\end{figure}

\subsubsection{RQ3: Does Dependency-Granularity Tree Awareness Improve Resource Efficiency?}
Comparing B6 (Full GranuWarm: $J = 0.07$, Mem = 280.49~MB, P95 = 43.24~ms) against monolithic B2 ($J = 1.05$, Mem = 5,888.00~MB) and B5 (Fuzzy+PSO without Granularity: $J = 1.64$, Mem = 2,794.73~MB, P95 = 15,015.65~ms) reveals the single most impactful architectural result: \textbf{a 95.7\% reduction in composite fitness cost over non-granular fuzzy-PSO baselines and a 95.2\% memory reduction over monolithic static prewarming}. By decoupling functions into shareable runtime and library bundles, GranuWarm warms common layers (Python 3.11, Node.js 20, NumPy, Pandas) that simultaneously accelerate multiple leaf functions within a compact 280~MB footprint.

\begin{figure}[htbp]
    \centering
    \resultplot{figure_5_live_validation.pdf}
    \caption{Live Kubernetes Validation: Empirical CDF of request latencies measured on a 3-node Kubernetes cluster running GranuWarm, confirming that simulation-derived P95 estimates ($\leq$~50~ms) hold under real container orchestration overhead.}
    \label{fig:live_validation}
\end{figure}

\subsection{Controller Runtime Overhead}

Table~\ref{tab:overhead} reports the execution overhead and resource utilization of GranuWarm's controller components.

\begin{table}[htbp]
\centering
\caption{Controller execution overhead and resource utilization.}
\label{tab:overhead}
\resizebox{\columnwidth}{!}{%
\begin{tabular}{lcccc}
\toprule
Component & Interval & CPU (\%) & Exec Time & Memory \\
\midrule
Prometheus Monitor & 10~s & 0.4\% & 12.3~ms & 28~MB \\
Fuzzy Engine (27 Rules) & 10~s & 0.8\% & 8.1~ms & 15~MB \\
Granularity Tree Projection & 10~s & 0.1\% & 1.2~ms & 5~MB \\
Constrained PSO ($S=30$) & 60~s & 3.5\% & 74.8~ms & 32~MB \\
Actuation Engine (Hysteresis) & 10~s & 0.2\% & 18.5~ms & 12~MB \\
\bottomrule
\end{tabular}%
}
\end{table}

The lightweight fast loop completes in under 22~ms, while the slow-loop PSO executes in 74.8~ms once every minute. Total controller memory consumption remains under 32~MB, confirming that the dual-timescale separation keeps operational overhead minimal.

\subsection{Statistical Significance}

We evaluate the statistical significance of GranuWarm's composite performance improvement using the non-parametric Wilcoxon signed-rank test across an expanded experimental evaluation matrix of $N = 10$ paired trials, spanning all five workload regimes (W1--W5) across two independent random seeds (Seeds 42 and 123). We test the standard two-tailed research hypothesis testing whether the median difference between baseline and GranuWarm composite costs is non-zero ($H_0: \text{Median}(J_{\text{Baseline}} - J_{\text{B6}}) = 0$ versus $H_1: \text{Median}(J_{\text{Baseline}} - J_{\text{B6}}) \neq 0$).

For $N = 10$ paired trials, the maximum possible sum of signed ranks is $\sum_{i=1}^{10} i = \frac{10 \times 11}{2} = 55$. In all comparisons against all five baseline configurations (B1 through B5), GranuWarm strictly dominates the baseline across every single trial ($J(\text{B6}) < J(\text{Baseline}) \; \forall k$), yielding a test statistic of $W = 0.0$. Under the null hypothesis, the exact two-tailed permutation test $p$-value for $W = 0$ with $N = 10$ is:
\begin{equation}
p = 2 \cdot \left(\frac{1}{2}\right)^{10} = \frac{2}{1024} = \frac{1}{512} \approx 0.00195
\end{equation}
which strictly satisfies two-tailed statistical significance at both $\alpha = 0.05$ and $\alpha = 0.01$ ($p < 0.002$). Asymptotic normal approximation with continuity correction ($\mu_W = \frac{N(N+1)}{4} = 27.5$, $\sigma_W = \sqrt{\frac{N(N+1)(2N+1)}{24}} = \sqrt{96.25} \approx 9.811$) yields:
\begin{equation}
z = \frac{W - \mu_W + 0.5}{\sigma_W} = \frac{0 - 27.5 + 0.5}{9.811} \approx -2.752
\end{equation}
corresponding to a two-tailed asymptotic $p = 2 \cdot \Phi(-2.752) \approx 0.0059$ ($p < 0.01$). Specifically:
\begin{itemize}[leftmargin=*]
    \item \textbf{B6 vs.\ B1 (No Prewarm):} $W = 0.0, \, p_{\text{exact}} = 0.00195$ (Statistically significant at $\alpha = 0.01$, $p < 0.01$).
    \item \textbf{B6 vs.\ B2 (Fixed Prewarm):} $W = 0.0, \, p_{\text{exact}} = 0.00195$ (Statistically significant at $\alpha = 0.01$, $p < 0.01$). GranuWarm decisively dominates static fixed prewarming across all 10 trials by eliminating the 5,888~MB monolithic memory footprint and preventing concurrency queuing during bursts.
    \item \textbf{B6 vs.\ B3 (Fuzzy-Only):} $W = 0.0, \, p_{\text{exact}} = 0.00195$ (Statistically significant at $\alpha = 0.01$, $p < 0.01$).
    \item \textbf{B6 vs.\ B4 (PSO-Only):} $W = 0.0, \, p_{\text{exact}} = 0.00195$ (Statistically significant at $\alpha = 0.01$, $p < 0.01$).
    \item \textbf{B6 vs.\ B5 (Fuzzy+PSO No Granularity):} $W = 0.0, \, p_{\text{exact}} = 0.00195$ (Statistically significant at $\alpha = 0.01$, $p < 0.01$).
\end{itemize}
These results confirm that GranuWarm's latency and resource improvements are statistically robust across diverse traffic regimes under rigorous two-tailed non-parametric testing.

%% ====================================================================
%% 7. LIMITATIONS AND THREATS TO VALIDITY
%% ====================================================================
\section{Limitations and Threats to Validity}

\textbf{Simulation-to-Reality Gap:} Our evaluation employs a calibrated discrete-event simulation environment matching empirical microbenchmark latencies. In live container platforms, physical layer sharing requires underlying system support (e.g., CRIU checkpoints, shared memory volumes, or zygote fork architectures like RainbowCake~\cite{rainbowcake2024} and Pagurus~\cite{pagurus2022}).

\textbf{Scalability to Large Fleets:} For fleets exceeding hundreds of functions, PSO complexity may increase. Hierarchical cluster partitioning or greedy pruning can bound optimization time.

\textbf{Workload Representativeness:} While our workloads cover five canonical patterns, real-world multi-tenant production traces may contain complex cross-function correlations.

%% ====================================================================
%% 8. CONCLUSION
%% ====================================================================
\section{Conclusion}

This paper presented GranuWarm, an adaptive closed-loop prewarming framework that mitigates serverless cold starts by unifying Mamdani fuzzy inference, dependency-granularity tree modeling, and constrained discrete PSO under a dual-timescale control loop. Evaluated across five canonical workload regimes with $N = 10$ paired trials, GranuWarm achieves an optimal Pareto tradeoff: it lowers P95 tail latency to 43.24~ms and SLA violations to 0.90\% under a compact 280.49~MB footprint, reducing composite fitness cost by 93.3\% compared to monolithic fixed prewarming (B2) and by 95.7\% compared to non-granular fuzzy-PSO baselines (B5), with two-tailed statistical significance ($p = 0.00195$) across all baselines. Future work will explore integrating adaptive neuro-fuzzy inference (ANFIS) and evaluating GranuWarm atop physical layer-caching container runtimes.

%% ====================================================================
%% REFERENCES
%% ====================================================================
\begin{thebibliography}{30}

\bibitem{jonas2019cloud}
E.~Jonas, J.~Schleier-Smith, V.~Sreekanti, C.-C.~Tsai, A.~Khandelwal, Q.~Pu, V.~Shankar, J.~Carreira, K.~Krauth, N.~Yadwadkar, J.~E.~Gonzalez, R.~A.~Popa, I.~Stoica, and D.~A.~Patterson, ``Cloud programming simplified: A Berkeley view on serverless computing,'' \emph{arXiv preprint arXiv:1902.03383}, 2019.

\bibitem{knative2024}
Knative Authors, ``Knative Serving Documentation,'' 2024. [Online]. Available: \url{https://knative.dev/docs/serving/}

\bibitem{openfaas2024}
A.~Ellis, ``OpenFaaS --- Serverless Functions Made Simple,'' 2024. [Online]. Available: \url{https://www.openfaas.com/}

\bibitem{openwhisk2024}
Apache Software Foundation, ``Apache OpenWhisk,'' 2024. [Online]. Available: \url{https://openwhisk.apache.org/}

\bibitem{mohan2019agile}
A.~Mohan, H.~Sane, K.~Doshi, S.~Edber, S.~Kompella, R.~Lakshman, and B.~Pal, ``Agile cold starts for scalable serverless,'' in \emph{Proc.\ 11th USENIX Workshop on Hot Topics in Cloud Computing (HotCloud)}, 2019.

\bibitem{wang2018peeking}
L.~Wang, M.~Li, Y.~Zhang, T.~Ristenpart, and M.~Swift, ``Peeking behind the curtains of serverless platforms,'' in \emph{Proc.\ USENIX Annual Technical Conference (ATC)}, 2018, pp.~133--146.

\bibitem{bermbach2020using}
D.~Bermbach, A.-S.~Karakaya, and S.~Buchholz, ``Using application knowledge to reduce cold starts in FaaS services,'' in \emph{Proc.\ ACM Symposium on Applied Computing (SAC)}, 2020, pp.~134--143.

\bibitem{shahrad2020serverless}
M.~Shahrad, R.~Fonseca, I.~Goiri, G.~Chaudhry, P.~Batum, J.~Cooke, E.~Laureano, C.~Tresness, M.~Russinovich, and R.~Bianchini, ``Serverless in the wild: Characterizing and optimizing the serverless workload at a large cloud provider,'' in \emph{Proc.\ USENIX Annual Technical Conference (ATC)}, 2020, pp.~205--218.

\bibitem{akkus2018sand}
I.~E.~Akkus, R.~Chen, I.~Rimac, M.~Stein, K.~Satzke, A.~Beck, P.~Aditya, and V.~Hilt, ``SAND: Towards high-performance serverless computing,'' in \emph{Proc.\ USENIX Annual Technical Conference (ATC)}, 2018, pp.~923--935.

\bibitem{oakes2018sock}
E.~Oakes, L.~Yang, D.~Zhou, K.~Houck, T.~Harter, A.~Arpaci-Dusseau, and R.~Arpaci-Dusseau, ``SOCK: Rapid task provisioning with serverless-optimized containers,'' in \emph{Proc.\ USENIX Annual Technical Conference (ATC)}, 2018, pp.~57--70.

\bibitem{suresh2021servermore}
A.~Suresh and A.~Gandhi, ``ServerMore: Opportunistic execution of serverless functions in the cloud,'' in \emph{Proc.\ ACM Symposium on Cloud Computing (SoCC)}, 2021, pp.~460--473.

\bibitem{shahrad2019architectural}
M.~Shahrad, J.~Balkind, and D.~Wentzlaff, ``Architectural implications of function-as-a-service computing,'' in \emph{Proc.\ IEEE/ACM International Symposium on Microarchitecture (MICRO)}, 2019, pp.~1063--1075.

\bibitem{daraghmeh2024optimizing}
M.~Daraghmeh, A.~Agarwal, and Y.~Jararweh, ``Optimizing serverless computing: A comparative analysis of multi-output regression models for predictive function invocations,'' \emph{Simulation Modelling Practice and Theory}, vol.~134, article 102925, 2024.

\bibitem{lin2021modeling}
P.~Lin and H.~Khazaei, ``Modeling and optimization of performance and cost of serverless applications,'' \emph{IEEE Transactions on Parallel and Distributed Systems}, vol.~32, no.~3, pp.~615--632, 2021.

\bibitem{rainbowcake2024}
Y.~Kang, P.~Sun, and X.~Shen, ``RainbowCake: Layer-wise container caching and sharing for serverless computing,'' in \emph{Proc.\ ACM International Conference on Architectural Support for Programming Languages and Operating Systems (ASPLOS)}, 2024, pp.~412--426.

\bibitem{pagurus2022}
Z.~Zheng, L.~Ma, and H.~Shen, ``Pagurus: Warm starter for serverless computing with container sharing,'' in \emph{Proc.\ USENIX Annual Technical Conference (ATC)}, 2022, pp.~685--700.

\bibitem{faascache2021}
A.~Fuerst and P.~Sharma, ``FaasCache: Keeping serverless computing clean with data-aware caching,'' in \emph{Proc.\ ACM ASPLOS}, 2021, pp.~586--599.

\bibitem{hanaforoosh2025granularity}
M.~Hanaforoosh et al., ``Granularity-tree-based cold-start scheduling for serverless functions,'' \emph{Future Generation Computer Systems}, vol.~164, article 107604, 2025.

\bibitem{buyya2023pso}
A.~S.~S.~M.~Buyya et al., ``Particle swarm optimization for task scheduling in cloud computing: A survey,'' \emph{Swarm and Evolutionary Computation}, vol.~76, article 101224, 2023.

\bibitem{zhan2022survey}
Z.~Zhan, X.~Shi, Z.~Lu, and H.~Lin, ``A survey on evolutionary computation for complex continuous optimization,'' \emph{Artificial Intelligence Review}, vol.~55, pp.~1--52, 2022.

\bibitem{mamdani1975experiment}
E.~H.~Mamdani and S.~Assilian, ``An experiment in linguistic synthesis with a fuzzy logic controller,'' \emph{International Journal of Man-Machine Studies}, vol.~7, no.~1, pp.~1--13, 1975.

\bibitem{jamshidi2014autonomic}
P.~Jamshidi, A.~Ahmad, and C.~Pahl, ``Autonomic resource provisioning for cloud-based software,'' in \emph{Proc.\ 9th International Symposium on Software Engineering for Adaptive and Self-Managing Systems (SEAMS)}, 2014, pp.~95--104.

\bibitem{kennedy1995pso}
J.~Kennedy and R.~Eberhart, ``Particle swarm optimization,'' in \emph{Proc.\ IEEE International Conference on Neural Networks}, vol.~4, 1995, pp.~1942--1948.

\bibitem{deb2000constraint}
K.~Deb, ``An efficient constraint handling method for genetic algorithms,'' \emph{Computer Methods in Applied Mechanics and Engineering}, vol.~186, no.~2--4, pp.~311--338, 2000.

\end{thebibliography}

\end{document}